% Figure: a pinned-pinned column under axial load P, with the exact first
% buckling mode (a half sine) and a one-term parabolic Rayleigh-Ritz trial
% shape laid over it for comparison. The two shapes nearly coincide, which
% is why the crude trial recovers the buckling load to better than a per cent.
\begin{figure}[htbp]
\centering
\begin{tikzpicture}[
  axis/.style={draw=engblack, thick},
  exact/.style={draw=engblue, very thick},
  trial/.style={draw=engred, very thick, dashed},
  loadarrow/.style={draw=engblack, -{Stealth}, thick},
  pin/.style={draw=engblack, fill=enggray!30},
  label/.style={font=\small},
]
% The undeformed column is the vertical segment from (0,0) to (0,6).
% We draw the buckled shapes as horizontal offsets w(s) for s in 0..6,
% amplitude 1.4 units at mid-height. Vertical s maps directly to y.
% Exact mode: w = A sin(pi s / 6).
\draw[axis] (0,0) -- (0,6);
\node[font=\scriptsize, enggray, anchor=east] at (-0.15,3) {undeformed};

% pin supports (triangles) at base and top
\fill[pin] (-0.22,-0.38) -- (0.22,-0.38) -- (0,0) -- cycle;
\fill[pin] (-0.22,6.38) -- (0.22,6.38) -- (0,6) -- cycle;
\node[font=\scriptsize, anchor=north] at (0,-0.55) {pinned};
\node[font=\scriptsize, anchor=south] at (0,6.7) {pinned};

% axial load arrows at top (pushing down) and base reaction
\draw[loadarrow] (0,8.0) -- (0,7.15);
\node[label, anchor=west] at (0.1,7.6) {$P$};

% exact first mode: half-sine, amplitude 1.4
\draw[exact] plot[domain=0:6, samples=80]
  ({1.4*sin(\x/6*180)}, {\x});
\node[engblue, label, anchor=west] at (1.5,4.5)
  {exact: $w=A\sin\dfrac{\pi s}{L}$};

% one-term parabolic trial: w = B s (L - s), normalised so peak = 1.4 at s=3.
% peak of s(L-s) at s=3 is 9, so B = 1.4/9 = 0.15556.
\draw[trial] plot[domain=0:6, samples=60]
  ({0.15556*\x*(6-\x)}, {\x});
\node[engred, label, anchor=west] at (1.0,1.6)
  {trial: $w=B\,s(L-s)$};

\end{tikzpicture}
\caption[Exact and trial buckling shapes for a pinned-pinned column]{The
first buckling mode of a pinned-pinned column (solid blue, a half-sine)
beside a one-term parabolic Rayleigh-Ritz trial shape (red dashed). Both
satisfy the geometric end conditions $w=0$ at each pin, and the two curves
nearly overlap to the eye. The closeness of the shapes is deceptive: the
parabola carries constant curvature where the true mode has curvature
peaking at mid-span and vanishing at the ends, and the buckling quotient
weights curvature heavily, so the one-term parabola over-estimates the
buckling load by about twenty per cent. A trial shape that also matches the
zero-moment end condition $w''=0$ at the pins brings the estimate within a
fraction of a per cent, the convergence the case study carries through.}
\label{fig:vol03ch05:column-buckling-modes}
\end{figure}
